\chapter{Kidney Infection (Pyelonephritis)}
\index{Kidney Infection (Pyelonephritis)}

\section*{Definition and Overview}
Pyelonephritis is the medical name for a kidney infection, an infection of one or both kidneys that occurs when bacteria, usually from the bladder, travel up the ureters into the kidney. It is a more serious infection than cystitis, affecting the whole body, with fever, chills, flank pain, and nausea. Pyelonephritis requires antibiotic treatment, often in hospital, and prompt treatment prevents complications such as kidney abscess, scarring, and sepsis. It is most common in women, and special care is needed in pregnancy, where it can harm both mother and baby. This chapter covers the causes, diagnosis, treatment, and prevention of kidney infection.

\section*{Epidemiology}
Pyelonephritis affects around 1 to 2 per cent of women each year and is less common in men. It is one of the most common reasons for hospital admission for infection, and it occurs in around 1 to 2 per cent of pregnancies, where it is the most common cause of serious non-obstetric infection. Risk is higher in women, in people with kidney stones, diabetes, urinary catheters, or obstruction, and after urinary procedures. Recurrence affects a significant minority, and repeated infections can cause kidney scarring, particularly in children with vesicoureteral reflux.

\section*{Aetiology and Pathophysiology}
Pyelonephritis begins when bacteria, most commonly Escherichia coli from the bowel, colonise the bladder and ascend the ureters into the kidney. Once in the kidney, the bacteria multiply in the renal tissue, triggering inflammation, swelling, and the release of toxins. Bacteria can enter the bloodstream, causing bacteraemia and the systemic features of the illness. Ascending infection is promoted by anything that impairs urine flow or allows reflux, including obstruction from stones or an enlarged prostate, pregnancy, which dilates the ureters, and reflux in children. In severe or untreated infection, pus can collect in abscesses, and the kidney can be permanently scarred, reducing function.

\section*{Clinical Features}
\begin{itemize}
    \item Fever, often with chills and rigors
    \item Flank or loin pain, usually on one side, and tenderness over the kidney
    \item Pain or burning with urination, frequency, and urgency
    \item Cloudy, smelly, or blood-stained urine
    \item Nausea, vomiting, and loss of appetite
    \item Fatigue, headache, and muscle aches
    \item Confusion in older people
    \item Fever, vomiting, and poor feeding in children
    \item Rapid onset, usually within hours to a day
\end{itemize}

\section*{Differential Diagnosis}
The diagnosis of pyelonephritis is supported by excluding other causes.
\begin{itemize}
    \item \textbf{Lower urinary tract infection:} cystitis without fever or flank pain
    \item \textbf{Kidney stones:} colicky pain, with infection if fever is present
    \item \textbf{Kidney abscess:} prolonged fever despite antibiotics
    \item \textbf{Abdominal conditions:} appendicitis, diverticulitis, and cholecystitis
    \item \textbf{Musculoskeletal pain:} back strain and disc disease
    \item \textbf{Sepsis from other sites:} pneumonia and other infections
    \item \textbf{Other renal disease:} glomerulonephritis and infarction
\end{itemize}

\section*{When to Seek Medical Care}
Fever with flank pain or a urinary infection that makes you feel generally unwell requires prompt assessment, and pregnant women, people with diabetes, and older people should seek care early. Anyone vomiting, unable to take antibiotics, or showing signs of sepsis needs urgent care.

\textbf{Contact your local emergency services for:}
\begin{itemize}
    \item High fever with shaking, confusion, or collapse, which may indicate sepsis
    \item Severe flank or abdominal pain
    \item Vomiting with an inability to keep fluids down
    \item Difficulty breathing or rapid deterioration
    \item Pregnancy with fever and flank pain
\end{itemize}

\section*{Diagnosis and Investigations}
\begin{itemize}
    \item \textbf{Urinalysis and culture:} confirming infection and the organism
    \item \textbf{Blood tests:} full blood count, inflammatory markers, kidney function, and blood cultures
    \item \textbf{Pregnancy test:} in women of reproductive age
    \item \textbf{Ultrasound or CT:} for stones, obstruction, and abscess
    \item \textbf{Assessment in children:} imaging for reflux and scarring
    \item \textbf{Follow-up:} repeat urine tests after treatment
\end{itemize}

\section*{Management}
\begin{itemize}
    \item \textbf{Antibiotics:} oral for mild infections, intravenous for severe illness and pregnancy
    \item \textbf{Hospital admission:} for sepsis, vomiting, pregnancy, or high-risk patients
    \item \textbf{Fluids and pain relief:} supportive care
    \item \textbf{Treatment of obstruction:} relieving stones or other blockage
    \item \textbf{Abscess drainage:} for collections of pus
    \item \textbf{Prophylactic antibiotics:} in some cases of recurrent infection
    \item \textbf{Follow-up:} confirming resolution and kidney recovery
    \item \textbf{Management in pregnancy:} hospital care and monitoring of the baby
\end{itemize}

\section*{Complications}
\begin{itemize}
    \item Sepsis and septic shock
    \item Kidney abscess and scarring
    \item Reduced kidney function and high blood pressure
    \item Recurrent infections
    \item Complications in pregnancy: preterm birth and low birth weight
    \item Kidney failure in severe or repeated disease
    \item Complications of treatment
\end{itemize}

\section*{Prognosis and Outcomes}
With prompt antibiotic treatment, the outlook for pyelonephritis is excellent, and most people improve within a few days and recover fully within weeks. Severe cases, and those in pregnancy, respond well to hospital care, and complications are treatable when recognised early. Managing underlying causes such as stones and reflux prevents recurrence, and children with reflux benefit from early detection and follow-up.

\section*{Prevention and Self-Care}
\begin{itemize}
    \item Seek treatment for bladder infections promptly
    \item Drink plenty of fluids
    \item Empty the bladder regularly and fully
    \item Complete antibiotic courses as prescribed
    \item Attend follow-up tests
    \item Manage diabetes and other risk factors
    \item Pregnant women should report urinary symptoms early
    \item Children with urinary infections should be followed up
    \item Seek urgent care for fever, severe pain, or vomiting
\end{itemize}

\section*{Sources and Further Reading}
The following reputable sources provide further information for healthcare students and patients seeking reliable, current advice about pyelonephritis.
\begin{itemize}
    \item Kidney Health Australia. \textit{Kidney infection information.} https://kidney.org.au/
    \item Healthdirect Australia. \textit{Pyelonephritis.} https://www.healthdirect.gov.au/
    \item Therapeutic Guidelines. \textit{Urinary tract infection in adults.} https://www.tg.org.au/
    \item Pregnancy, Birth and Baby. \textit{Kidney infections in pregnancy.} https://www.pregnancybirthbaby.org.au/
    \item NHS. \textit{Kidney infection.} https://www.nhs.uk/conditions/kidney-infection/
    \item Mayo Clinic. \textit{Kidney infection.} https://www.mayoclinic.org/
\end{itemize}
